\chapter{Mathematical and Cryptological Synthesis}
\label{ch:mathematical-cryptological-synthesis}

This chapter gives the central mathematical and cryptological reading of the
machine-checked \haetae{} security proof.  It also repairs an important weakness
of a purely catalog-style dissertation: the final theorem must not be hidden
behind a long list of EasyCrypt names.  A cryptologist needs to see the security
experiment, the adversary model, the explicit loss ledger, the residual trusted
base, and the precise sense in which the EasyCrypt proof is non-vacuous.

The checked development proves a conditional theorem.  It verifies the game
hops, oracle-programming arguments, sampler-distribution arguments, budgeted
active-signing lifting, and final reduction composition.  It does not prove the
underlying lattice assumptions and it does not by itself prove that every
implementation realizes the formal model.  The dissertation should therefore be
read as a rigorous mechanization of the provable-security reduction for the
formal \haetae{} model.

\section{The formal security object}
\label{sec:synthesis-formal-object-rewritten}

Let \(A\) be an adversary with access to a hash oracle \(H\) and a signing
oracle \(O\).  The random-oracle EUF-CMA experiment samples
\[
  (pk,sk) \leftarrow \mathsf{KeyGen}(),
\]
answers signing queries by \(\mathsf{Sign}^{H}(sk,\cdot)\), answers hash queries
by lazy sampling, and receives a final output
\[
  (m^\star,ctx^\star,\sigma^\star)\leftarrow A^{H,O}(pk).
\]
The success event is
\[
  \mathsf{Forge}
  =
  \left[
    \mathsf{Verify}^{H}(pk,m^\star,ctx^\star,\sigma^\star)=1
    \wedge
    (m^\star,ctx^\star)\notin Q_S
  \right],
\]
where \(Q_S\) is the signing-query log.  The advantage is
\[
  \Adv_{\haetae,\rom}^{\eufcma}(A)=\Pr[\mathsf{Forge}].
\]

In EasyCrypt this probability is represented by a moduleized experiment in
\ec{Sig\_ROM.ec}, the scheme interface in \ec{HAETAE\_Scheme.ec}, and the hop
layers in \ec{HAETAE\_GameHops.ec} and \ec{HAETAE\_HopGames.ec}.  The theorem
proved by the checker has the mathematical shape
\[
  \Adv_{\haetae,\rom}^{\eufcma}(A)
  \le
  \ec{haetae\_euf\_counted\_rom\_bound}.
\]
The identifier on the right is not a black-box placeholder; it is the concrete
EasyCrypt expression assembled and normalized in \ec{HAETAE\_Reductions.ec}.

\section{Expanded advantage ledger}
\label{sec:synthesis-expanded-ledger}

For a dissertation intended for cryptographers, the bound should be displayed as
a ledger.  The following theorem gives the paper-level version of the checked
statement.

\begin{theorem}[Ledger form of the checked \haetae{} security theorem]
\label{thm:ledger-form-haetae-security}
For every adversary \(A\) satisfying the declared hash-query and signing-query
budgets, there is a reduction adversary \(B_A\) such that
\[
\begin{aligned}
  \Adv_{\haetae,\rom}^{\eufcma}(A)
  \le{}&
  L_{\mathsf{MLWE}}
  + L_{\mathsf{ModSIS}}
  + L_{\mathsf{bimodal\mbox{-}MSIS}}
  + L_{\mathsf{ROM\mbox{-}coll}}
  + L_{\mathsf{ROM\mbox{-}pre}} \\
  &+ L_{\mathsf{ROM\mbox{-}me}}
  + L_{\mathsf{sample}}
  + L_{\mathsf{rej}}
  + L_{\mathsf{extract}}
  + L_{\mathsf{budget}} .
\end{aligned}
\]
The EasyCrypt theorem states this same inequality with the right-hand side
encoded as \ec{haetae\_euf\_counted\_rom\_bound}; the equivalence between the
ledger notation and the checked expression is witnessed by
\ec{haetae\_euf\_boundE}, \ec{haetae\_euf\_bound\_groupedE}, and
\ec{haetae\_euf\_counted\_rom\_bound\_groupedE}.
\end{theorem}

\begin{longtable}{p{0.20\textwidth}p{0.36\textwidth}p{0.34\textwidth}}
\caption{Cryptological loss ledger for the checked security theorem.}
\label{tab:cryptological-loss-ledger}\\
\toprule
\textbf{Paper term} & \textbf{Meaning} & \textbf{EasyCrypt evidence} \\
\midrule
\endfirsthead
\toprule
\textbf{Paper term} & \textbf{Meaning} & \textbf{EasyCrypt evidence} \\
\midrule
\endhead
\(L_{\mathsf{MLWE}}\) & Assumed advantage term for the MLWE-style hardness game
used by the reduction layer. & \ec{mlwe\_hardness\_term} in
\ec{HAETAE\_Reductions.ec} and the assumption interfaces in
\ec{HAETAE\_Assumptions.ec}. \\
\midrule
\(L_{\mathsf{ModSIS}}\) & Assumed Module-SIS advantage term. &
\ec{module\_sis\_hardness\_term} and the Module-SIS game interfaces. \\
\midrule
\(L_{\mathsf{bimodal\mbox{-}MSIS}}\) & Bimodal self-target MSIS term produced by
forgery extraction and special-soundness reasoning. &
\ec{bimodal\_selftarget\_msis\_hardness\_term},
\ec{BimodalSelfTargetMSIS\_Game}, and reduction modules in
\ec{HAETAE\_HopGames.ec}. \\
\midrule
\(L_{\mathsf{ROM\mbox{-}coll}}\) & Loss from collisions among programmed or
observed random-oracle sites. & \ec{counted\_rom\_collision\_term} and grouped
ROM-programming lemmas. \\
\midrule
\(L_{\mathsf{ROM\mbox{-}pre}}\) & Loss that a fresh sampler-programming query was
already asked by the adversary. & \ec{counted\_rom\_prequery\_term},
\ec{sampler\_bad\_prequery}, and logged adversary hash-query lemmas. \\
\midrule
\(L_{\mathsf{ROM\mbox{-}me}}\) & Min-entropy or support-size loss needed to turn
fresh coins into a concrete prequery probability. &
\ec{counted\_rom\_min\_entropy\_term} and spread/support lemmas. \\
\midrule
\(L_{\mathsf{sample}}\) & Statistical distance between attempt, paper, and exact
sampling distributions. & Structural-to-exact and exact-hyperball sampler
lemmas in \ec{HAETAE\_Distributions.ec}, \ec{HAETAE\_Rejection.ec}, and
\ec{HAETAE\_HopGames.ec}. \\
\midrule
\(L_{\mathsf{rej}}\) & Loss paid for rejection-sampling conditioning and support
membership. & Rejection predicates and loss lemmas in \ec{HAETAE\_Rejection.ec}. \\
\midrule
\(L_{\mathsf{extract}}\) & Loss in converting a final forgery or forked transcript
into a hardness-game solution. & Transcript-validity, special-soundness, and
bimodal extraction lemmas in \ec{HAETAE\_Transcript.ec} and
\ec{HAETAE\_HopGames.ec}. \\
\midrule
\(L_{\mathsf{budget}}\) & Multiplicative loss from lifting one-call bounds through
adaptive hash and signing oracle interaction. & Budgeted wrapper lemmas over
\ec{ROMInternalTranscriptBudgetedPaperSimAsNMA.O.sign}. \\
\bottomrule
\end{longtable}

This ledger is intentionally more informative than the raw EasyCrypt expression.
The formal expression is the object checked by EasyCrypt; the ledger is the
cryptological interpretation of its summands.

\section{Denotational reading of the proof scripts}
\label{sec:synthesis-denotational-rewritten}

Each EasyCrypt procedure denotes a subprobability kernel.  If \(S\) is the state
space and \(R\) is the return space of a procedure \(c\), then
\[
  \llbracket c\rrbracket:S\to\mathcal{D}(S\times R).
\]
Losslessness lemmas assert total mass one.  Hoare judgments assert almost-sure
postcondition preservation:
\[
  \{P\}\;c\;\{Q\}
  \quad\Longleftrightarrow\quad
  \forall s.\;P(s)\Rightarrow
  \Pr_{(s',r)\leftarrow\llbracket c\rrbracket(s)}[Q(s',r)]=1.
\]
Equivalence judgments are relational kernel statements.  A judgment
\[
  \mathbf{equiv}[c_1\sim c_2:P\Rightarrow Q]
\]
asserts that, for related initial states, the two output distributions admit a
coupling under which \(Q\) holds almost surely.  This is the mathematical reason
a game-hop proof in EasyCrypt can be read as a conventional cryptographic hybrid
argument.

\begin{proposition}[Forms of checked game transitions]
\label{prop:checked-game-transition-forms}
Every security-relevant transition in the proof has one of the following forms:
\[
  \Pr[G_i:E_i]=\Pr[G_{i+1}:E_{i+1}],
\]
\[
  \Pr[G_i:E_i]\le \Pr[G_{i+1}:E_{i+1}]+\Pr[G_i:\mathsf{Bad}_i],
\]
\[
  \Pr[G_i:E_i]\le B_i(A).
\]
The first is a perfect coupling, the second is a fundamental-lemma transition,
and the third is a direct probability or reduction bound.
\end{proposition}

\section{Adversary budgets and reduction cost}
\label{sec:synthesis-reduction-cost}

The adversary model is budgeted.  Let
\[
  q_H=|Q_H|,
  \qquad
  q_S=|Q_S|,
\]
where \(Q_H\) is the hash-query log and \(Q_S\) is the signing-query log.  The
proof must preserve these inequalities under every oracle call.  A one-call loss
\(\epsilon_{\mathsf{sign}}\) therefore lifts to at most
\(q_S\epsilon_{\mathsf{sign}}\), and a one-hash-call loss
\(\epsilon_{\mathsf{hash}}\) lifts to at most
\(q_H\epsilon_{\mathsf{hash}}\).

A complete cryptographic statement should also record reduction cost.  The
current EasyCrypt proof is strongest on probability accounting and oracle-call
accounting; it is less explicit about wall-clock or circuit-size overhead.  The
paper-level reduction-cost statement should be maintained as the following
ledger:
\[
  \operatorname{Time}(B_A)
  \le
  \operatorname{Time}(A)
  + T_{\mathsf{sim}}(q_H,q_S)
  + T_{\mathsf{extract}}(q_H,q_S),
\]
where \(T_{\mathsf{sim}}\) is the cost of simulated signing, lazy-ROM table
maintenance, and sampler instrumentation, and \(T_{\mathsf{extract}}\) is the
cost of the extraction/forking layer.  If future EasyCrypt files formalize time
costs, this ledger should be replaced by checked cost bounds; until then it is a
cryptographic interpretation obligation, not a machine-checked theorem.

\section{Non-vacuity of the clean-conditioned proof}
\label{sec:synthesis-non-vacuity-rewritten}

A machine-checked proof can be formally accepted but cryptologically vacuous if
the hypotheses of the main game-hop theorem are never reachable.  The active
signing part of the \haetae{} proof avoids this only if the clean condition is
initialized, preserved, and strong enough to trigger the next transition.

Let
\[
  C_i =
  \neg\mathsf{sampler\_bad\_prequery}
  \wedge
  \mathsf{sampler\_rom\_covered}
  \wedge
  pk_{\mathsf{current}}=pk
  \wedge
  \mathsf{ROMConsistent}_i.
\]
The non-vacuity obligation is the triple
\[
  \mathsf{Init}\Rightarrow C_0,
\]
\[
  C_i\wedge\neg\mathsf{Bad}_i\Rightarrow C_{i+1},
\]
\[
  C_i\Rightarrow
  \text{the clean-conditioned lifting theorem is applicable at step }i.
\]
The EasyCrypt evidence consists of the sampler-ROM coverage lemmas, the concrete
lazy-ROM freshness laws, and the active \ec{O.sign} lifting lemmas, especially
\ec{concrete\_budgeted\_self\_logged\_signature\_clean\_event\_loss\_surface}
and
\ec{concrete\_budgeted\_o\_sign\_clean\_one\_call\_loss\_from\_self\_logged\_surface}.

\begin{lemma}[Paper-level non-vacuity criterion]
\label{lem:paper-level-non-vacuity}
The clean-conditioned lifting is cryptologically meaningful only if the
initialization, preservation, and applicability obligations above are all
available for the concrete \haetae{} lazy ROM and the concrete active signing
oracle.
\end{lemma}

\begin{proof}
If initialization fails, the clean-conditioned theorem is never entered.  If
preservation fails, adaptive signing may leave the clean region before the game
hop is applied.  If applicability fails, the later sampler and ROM-programming
lemmas cannot be instantiated.  Hence all three obligations are necessary.  The
formal development addresses them through concrete logged-oracle and self-logged
signing instrumentation rather than by postulating a clean state abstractly.
\end{proof}

\section{Trusted base and residual obligations}
\label{sec:synthesis-trusted-base}

The following table is the paper's trust-boundary audit.  It states what is
checked, what is assumed, and what should be inspected when claiming security for
an implementation rather than for the formal model.

\begin{longtable}{p{0.23\textwidth}p{0.31\textwidth}p{0.36\textwidth}}
\caption{Trusted base and residual obligations.}
\label{tab:trusted-base-ledger}\\
\toprule
\textbf{Boundary item} & \textbf{Status} & \textbf{Required evidence} \\
\midrule
\endfirsthead
\toprule
\textbf{Boundary item} & \textbf{Status} & \textbf{Required evidence} \\
\midrule
\endhead
EasyCrypt logic and libraries & Trusted deductive base. & Versioned EasyCrypt
installation, standard library assumptions, and reproducible compile log. \\
\midrule
Hardness assumptions & Assumed, not proved. & Clear mapping from
\ec{HAETAE\_Assumptions.ec} games and terms to MLWE, Module-SIS, and bimodal
self-target MSIS assumptions. \\
\midrule
Parameter constants & Checked internally once declared. & Audit that
\ec{HAETAE\_Params.ec} matches the intended \haetae{} parameter set and security
level. \\
\midrule
Formal scheme model & Checked as a model. & Model/spec correspondence showing
that signing, verification, transcript formation, rejection sampling, and hash
queries match the HAETAE specification. \\
\midrule
Implementation fidelity & Outside this proof. & Separate refinement or
functional-correctness proof connecting concrete code to the EasyCrypt model. \\
\midrule
Manifest completeness & Procedural trust boundary. & Audited
\ec{provable-security/proof-files.txt}, compile script, and policy for adding or
removing proof files. \\
\midrule
Random-oracle abstraction & Idealized assumption. & The theorem is in the ROM;
concrete SHAKE/FIPS202 instantiation requires an additional cryptographic
argument. \\
\bottomrule
\end{longtable}

This table should remain in the main paper.  It prevents the phrase
``machine-checked security proof'' from being read as a claim that all external
assumptions and implementation links have also been proved.

\section{What FIPS202 and Keccak support do not prove}
\label{sec:synthesis-fips-boundary}

The files \ec{HAETAE\_FIPS202.ec} and \ec{HAETAE\_Keccak1600.ec} support the
modeled hash interface and concrete hash-domain objects used by the formalization.
They do not turn the random-oracle theorem into a standard-model proof for
SHAKE.  Likewise, the excluded generated KAT file is useful implementation
validation evidence, but it is not a game-hop theorem, a ROM-programming lemma,
or a hardness reduction.

Thus the correct statement is:
\[
  \text{the proof is a machine-checked ROM proof for the formal model,}
\]
not:
\[
  \text{the proof establishes standard-model security of the concrete hash
  implementation.}
\]
This distinction should be explicit wherever the dissertation discusses FIPS202
support.

\section{Why the generated catalog is supplemental}
\label{sec:synthesis-generated-catalog}

The exhaustive declaration catalog is valuable for traceability, but it is not a
substitute for mathematical exposition.  A generated line saying that a lemma
means ``the proposition stated by the lemma'' is formally true but not persuasive
to a cryptologist.  The main dissertation therefore keeps curated case studies
in the appendices and moves the exhaustive generated correspondence to
\ec{dissertation/etc/}.  The archived material remains available for audit, but
the paper's argument is now carried by theorem statements, ledgers, and selected
proof-object analyses.

\section{Summary}
\label{sec:synthesis-summary-rewritten}

The strengthened dissertation claim is now precise.  The EasyCrypt development
proves that a budget-respecting adversary against the formal random-oracle
\haetae{} signature model has success probability at most the explicit counted
bound \ec{haetae\_euf\_counted\_rom\_bound}.  The cryptological content of that
bound is the ledger of hardness, ROM, sampler, rejection, extraction, and
budgeted active-signing losses.  The remaining obligations are not hidden; they
are listed as trusted-base and model-fidelity requirements.
